% !TeX root = ../../surface-lagrangian-torus-fibration.tex
% 2026-11-22 / v1.8: 主定理の統合証明と許容中間尺度の帰結．
\section{Actual Lagrangian fibration による負 KE 極限の実現}
\label{ke:global-metric-realization-section}

\begin{theorem}[全固定次数における大域的な計量実現]
\label{ke:global-metric-realization}
定理\ref{ke:global-realization}の固定次数 $d\ge5$ と退化族に対し，
節\ref{ke:cut-selection-order}の順序で有限段階のデータを選ぶと，
同定理の二条件を同時に満たす一つの固定木系 $T$ と actual Lagrangian fibration
$\pi_t:(X_t,\omega_t)\to B_d(T)$ が得られる．特に，同じ全射
$f_t=p_T\pi_t$ に対して
\begin{equation}\label{ke:eq:global-metric-realization}
 \operatorname{dis}(f_t):=
 \sup_{x,y\in X_t}
 \left|R_t^{-1}d_{g_t}(x,y)
       -d_\square(f_t(x),f_t(y))\right|\longrightarrow0
\end{equation}
となる．距離は全曲面の ambient 距離であり，量化と規格化は
定理\ref{ke:global-realization}のものを用いる．固定 stage の明示的な評価は
式\eqref{ke:eq:global-realization-explicit-cut-budget}で与える．
\end{theorem}

\begin{proof}
\emph{構成と同じ写像の選択．}
定理\ref{ke:actual-global-fibration}により，固定 stage $j$ の後で
$|t|$ を十分小さくすれば，全 $X_t$ 上に actual fibration を得る．
全被覆，連結 fibers，actual Lagrangian 性はこの定理で証明されている．
補題\ref{ke:fixed-marked-tree-exhaustion}の共通選択により，
木系の組合せ型，labelled ends，および source markings は stage と $t$ に依存せず固定できる．
ここでは各固定二重線の一つの embedded pants foliation から構成し，
動く source を小円板内で，end germs をその外の annuli で相対的に調整している．
Normal circle の向きは共通にするが，全 torus cycles の parameter-monodromy を
自明化する主張ではない．
命題\ref{ke:global-base-parameterization}を用いて，その同じ fiber partition の
基底を $B_d(T)$ と parameter 化し，$f_{j,t}=p_T\pi_{j,t}$ とする．
これは距離比較のために別の fiber partition を採用する操作ではない．
補題\ref{ke:anchor-trimmed-neck-domain}により，core と変更 collar の外で使う
保持 neck は共通 profile の定義域 $|\widehat\mu|\le\delta_j$ に入る．
補題\ref{ke:same-moment-corner-reference}は，その同じ moment の旧 corner 辺への
接続を与える．下界の値比較と上界の portals は，この同じ固定残余基底と
canonical moment を用いている．

\emph{固定 stage の両側評価．}
定理\ref{ke:global-ambient-upper}の記号で
\[
 U_{j,t}=M_dE_{j,t}+C(M_d+1)\eta_{j,t}
          +(M_d+1)H_{j,t}/R_t,
 \qquad M_d=3\binom{\min\{d,6\}}3\le60
\]
とする．下界の誤差を
\begin{equation}\label{ke:eq:global-realization-lower-budget}
 V_{j,t}=C_d\left\{\varepsilon_j+b_j^{-1}+K_j^{-2}
 +\frac{1+\log\mathfrak m_j+\log K_j+\log(N_j+1)}{R_t}\right\}
\end{equation}
と書く．定理\ref{ke:global-ambient-lower}と
定理\ref{ke:global-ambient-upper}は，\emph{同じ} $f_{j,t}$ に対し
全 $x,y\in X_t$ で
\[
 d_\square(f_{j,t}(x),f_{j,t}(y))-V_{j,t}
 \le R_t^{-1}d_{g_t}(x,y)
 \le d_\square(f_{j,t}(x),f_{j,t}(y))+U_{j,t}
\]
を与える．上界では全内部木の共通 moment，source の全 moment 退出，
各 component compact の費用，および高々 $M_d$ 本の square gallery を使う．
下界では global clipped-depth map の長さ評価を全ての経路に直接積分する．
従って局所 protected 距離を ambient 距離に読み替える段階も，
経路の領域出入りの回数を有限と仮定する段階もない．
上界の式\eqref{ke:eq:global-upper-explicit-depth-error}と合わせると，
許容 cut の下で固定族の定数 $C_d$ により
\begin{equation}\label{ke:eq:global-realization-explicit-cut-budget}
 \operatorname{dis}(f_{j,t})\le C_d\left\{
 \varepsilon_j+b_j^{-1}+K_j^{-2}
 +\frac{1+\log R_t+\log\mathfrak m_j+\log K_j+\log(N_j+1)}{R_t}
 \right\}.
\end{equation}
ここには独立の stage-dependent な乗法係数を残していない．
Stage の定数は，cut の下限と最後の parameter 閾値へ入っている．

\emph{パラメータの選択順序と対角化．}
固定領域上で既に証明した有限個の誤差と費用を一つの遅い段階選択にまとめる原理は，
関連する局所解析原稿の補題\cite[\texttt{pair:lem:diagonal-fixed-end}]{KELocalCurrent}に明記した．
以下では本稿固有の cut の選択順序を確認する．この対角化自体は，
前段で別に証明した moving buffer，source 退出，scalar 下界を供給するものではない．
まず $b_j\to\infty$，$\varepsilon_j\downarrow0$ を選び，旧 corner の cut を
$m_j=6b_j$ とする．各 stage の residual domains とその接続定数
$C_j^{\rm neck}$ を先に固定する．次に component cut $N_j$ を，既存の下限および
\eqref{ke:eq:neck-constant-depth-choice}を満たすよう選び，core と相対 collars を作る．
この順序で $H_{j,t}=C(1+\log(N_j+1))$ となり，
上界には\eqref{ke:eq:global-upper-explicit-depth-error}が使える．
定理\ref{ke:sharp-global-one-active}と
命題\ref{ke:global-depth-base-comparison}の順序に従い，
Scalar 障壁専用の補助 cut $\mathfrak n_j$ とその境界定数を決めてから，
新しい lower cut $\mathfrak m_j$ と upper cut factor $K_j\to\infty$ を選ぶ．
$\mathfrak n_j$ は既存 neck/corner 構成の $n_j$ とは独立の記号である．
この $\mathfrak m_j$ は旧 corner cut $m_j$ と同じ量ではない．
構成・比較の有限個の条件に加えて
\begin{equation}\label{ke:eq:global-realization-diagonal}
 R_t\ge j\{1+\log\mathfrak m_j+\log K_j+\log(N_j+1)\}
\end{equation}
を満たす閾値 $r_j>0$ を取る．$r_{j+1}<r_j/2$ としてよく，
$0<|t|<r_j$ なら stage $j$ の全条件が成り立つようにする．
$j(t)=\max\{j:|t|<r_j\}$ と置けば $j(t)\to\infty$ であり，
\[
 \frac{\log(m_{j(t)}+1)+\log R_t}{R_t}\to0,
 \quad \frac{1+\log(N_{j(t)}+1)}{R_t}\to0.
\]
ここでは $N_j\ge m_j$ と，$\log R_t/R_t\to0$ を用いた．
従って $U_{j(t),t},V_{j(t),t}\to0$ となる．
$\pi_t=\pi_{j(t),t}$ と選び，両側評価から
$\operatorname{dis}(f_t)\le\max\{U_{j(t),t},V_{j(t),t}\}\to0$ を得る．
この定理は $t$ に関する fibration 族の連続性・滑らかさや，
あらかじめ指定した速さの対角化を結論に含めない．
明示的な cut budget は，$C_j^{\rm neck}$ の増大率，scalar 障壁に必要な
$\mathfrak m_j,K_j$ の成長，又は fixed-compact convergence の有効な閾値を
与えるものではない．従って $R_t$ だけによる最適な全域収束速度の証明とは区別する．
\end{proof}

\begin{lemma}[全域の歪みから同半径の閉球への移送]
\label{ke:distortion-to-pointed-balls}
Compact geodesic spaces $(X,d_X)$，$(Y,d_Y)$ と，全射 $F:X\to Y$ があり，
全点対に対する distortion が $\delta\ge0$ 以下とする．
任意の $x\in X$ と $A>0$ に対し，ambient 距離を制限した二つの閉球
\[
 (\overline B_X(x,A),x),\qquad
 (\overline B_Y(F(x),A),F(x))
\]
の間には，基点対を含む，distortion が $3\delta$ 以下の対応がある．
従って基点付き GH 距離は $3\delta/2$ 以下である．定数は $x,A$ と両空間に依存しない．
\end{lemma}
\begin{proof}
$z\in\overline B_X(x,A)$ に対し，$d_Y(F(x),F(z))\le A+\delta$ である．
必要なら $F(x)$ から $F(z)$ への最短路上で距離 $\delta$ 以下だけ戻し，
半径 $A$ の球内の点 $y_z$ を選び，対 $(z,y_z)$ を取る．基点対 $(x,F(x))$ も取る．
球内の任意の $y$ に対し，全射性から $F(w)=y$ となる $w$ を選ぶと
$d_X(x,w)\le A+\delta$ である．Source の最短路上で必要なだけ戻せば，
$d_X(z,w)\le\delta$ となる $z\in\overline B_X(x,A)$ を取れる．
この対 $(z,y)$ も加えると，両閉球へ全射な対応を得る．
この対応の各対 $(z,y)$ には，元の全域写像の点 $w$ で
\[
 d_X(z,w)+d_Y(y,F(w))\le\delta
\]
となるものが付随する．二つの対に対してこの評価と $F$ の distortion を加えれば，
対応の distortion は $3\delta$ 以下である．基点対を含む対応から通常の共通距離を
構成すると，両球の Hausdorff 距離と基点間距離は $3\delta/2$ 以下となる．
これらの選択の連続性は必要ない．
\end{proof}

\begin{corollary}[同じ射影で制御される中間尺度の全移動基点極限]
\label{ke:intermediate-moving-classification}
定理\ref{ke:global-metric-realization}の同じ写像 $f_t$ と
$\epsilon_t=\operatorname{dis}(f_t)$ を用い，正数 $a_t$ が
\[
 a_t\longrightarrow\infty,\qquad a_t/R_t\longrightarrow0,
 \qquad R_t\epsilon_t/a_t\longrightarrow0
\]
を満たすとする．任意の $t_\nu\to0$ と $x_\nu\in X_{t_\nu}$ に対し，
$(X_{t_\nu},a_{t_\nu}^{-1}d_{g_{t_\nu}},x_\nu)$ は部分列を取れば pointed GH 収束する．
その全ての可能な極限は，命題\ref{topo:cubical-moving-blowups}の四種類
\[
 (C_{2,d-1},z),\quad
 (\mathbb R\times\mathcal Y_{d-2},(0,z)),\quad
 (\mathcal K_{3\pi/2},z),\quad
 (\mathbb R^2,0)
\]
であり，この条件を満たす任意の尺度と退化列に対して，各型の任意の有限基点を
実現する $x_\nu$ を選べる．これは metric spaces の分類であり，
追加の labels 又は tensor 収束を結論に含めない．
\end{corollary}
\begin{proof}
$\lambda_t=R_t/a_t\to\infty$，$\delta_t=\lambda_t\epsilon_t\to0$ と置く．
同じ全射 $f_t$ は $(X_t,a_t^{-1}d_{g_t})$ と
$(\mathcal C_d,\lambda_t d_\square)$ の全点対距離を $\delta_t$ 以下の誤差で比較する．
両者は compact geodesic spaces なので，前補題により，任意の基点と固定半径の
閉球を $O(\delta_t)$ の誤差で比較できる．命題\ref{topo:cubical-moving-blowups}を
$q_\nu=f_{t_\nu}(x_\nu)$ に適用すれば，部分列の存在と分類が従う．
逆に同命題で各型を実現する $q_\nu$ を選び，全射性によって
$x_\nu\in f_{t_\nu}^{-1}(q_\nu)$ を取れば，その任意の選択が同じ極限を実現する．
必要なのは指定した尺度で $R_t\epsilon_t/a_t\to0$ となることであり，
この条件を除いた任意の中間尺度に関する分類ではない．
\end{proof}

\begin{corollary}[選択した中間尺度での成分錐]
\label{ke:intermediate-component-cone}
定理\ref{ke:global-metric-realization}の同じ $f_t$ に対し
$\epsilon_t=\operatorname{dis}(f_t)\to0$ と置く．成分 label $i$ を固定し，
任意に $x_{i,t}\in f_t^{-1}(e_i)$ を選ぶ．正数 $a_t$ が
\begin{equation}\label{ke:eq:intermediate-component-scale}
 a_t\longrightarrow\infty,\qquad
 \frac{a_t}{R_t}\longrightarrow0,\qquad
 \frac{R_t\epsilon_t}{a_t}\longrightarrow0
\end{equation}
を満たすなら，pointed Gromov--Hausdorff の意味で
\[
 (X_t,a_t^{-1}d_{g_t},x_{i,t})
 \longrightarrow (C_{2,d-1},d_C,0).
\]
このような尺度は常に存在し，例えば
\begin{equation}\label{ke:eq:intermediate-component-choice}
 a_t=R_t\sqrt{\epsilon_t+R_t^{-1}}
\end{equation}
と取れる．極限の錐は命題\ref{topo:component-cone-germ}のものである．

より一般に，非空集合 $A_t\subset X_t$ に対し
\[
 \eta_t=\sup_{y\in A_t}d_\square(f_t(y),e_i)\longrightarrow0
\]
とする．式\eqref{ke:eq:intermediate-component-scale}の第三条件を
\begin{equation}\label{ke:eq:intermediate-basepoint-budget}
 \frac{R_t(\epsilon_t+\eta_t)}{a_t}\longrightarrow0
\end{equation}
へ強めれば，任意の $y_t\in A_t$ を基点として同じ pointed GH 収束が成り立つ．
各固定半径の閉球に対する近似誤差は $y_t$ の選択に一様である．この場合には
\[
 a_t=R_t\sqrt{\epsilon_t+\eta_t+R_t^{-1}}
\]
と選べる．特に基点を成分頂点の特異 fiber 上へ限定する必要はない．
\end{corollary}
\begin{proof}
$\lambda_t=R_t/a_t$，$\delta_t=\lambda_t\epsilon_t$ と置くと，
$\lambda_t\to\infty$，$\delta_t\to0$ である．以下，閉球には各 ambient 距離の制限を入れる．
全射 $f_t$ は基点を $e_i$ へ送り，
$(X_t,a_t^{-1}d_{g_t})$ と $(\mathcal C_d,\lambda_t d_\square)$ の
距離の差を全点対について $\delta_t$ 以下にする．
固定した $A>0$ に対し，十分小さい $t$ では $(A+2\delta_t)/\lambda_t<1/2$ なので，
命題\ref{topo:component-cone-germ}により target のその半径の球は，
錐の半径 $A+2\delta_t$ の球と等長である．
補題\ref{ke:distortion-to-pointed-balls}により source と target の同半径 $A$ の閉球を
基点を保って GH 誤差 $3\delta_t/2$ 以下で比較できる．この等長写像と合成して
錐との比較を得るので，pointed GH 収束が従う．
式\eqref{ke:eq:intermediate-component-choice}は $a_t\ge\sqrt{R_t}$，
$a_t/R_t\to0$，$R_t\epsilon_t/a_t\le\sqrt{\epsilon_t}\to0$ を満たす．

一般の基点についても，同じ $x_{i,t}\in f_t^{-1}(e_i)$ を一つ取る．
全域 distortion 評価から，全 $y\in A_t$ に対して
\[
 a_t^{-1}d_{g_t}(y,x_{i,t})
 \le\frac{R_t}{a_t}(\eta_t+\epsilon_t)=:s_t\longrightarrow0
\]
である．同じ geodesic 空間の二つの中心が距離 $s_t$ 以下なら，その半径 $A$ の
閉球の Hausdorff 距離も $s_t$ 以下である．実際，一方の球の点を他方の中心から
最短路で結び，必要な場合だけ半径 $A$ まで戻せば，移動距離は高々 $s_t$ である．
逆方向も同じであり，中心間の距離も $s_t$ 以下だから，基点を含む近似を得る．
これを $x_{i,t}$ を中心とする上の近似と合成すれば，誤差は
$O(\delta_t+s_t)$ となり，全ての $y\in A_t$ に一様な結論が従う．
最後の尺度の選択では $a_t\ge\sqrt{R_t}$ であり，
\[
 \frac{a_t}{R_t}\to0,\qquad
 \frac{R_t(\epsilon_t+\eta_t)}{a_t}
 \le\sqrt{\epsilon_t+\eta_t}\to0
\]
となる．
\end{proof}

\paragraph{Component core と変更 collar 内の基点．}
同じ $f_t$ の構成で共通 anchor によって定めた component $i$ の core domain と
変更 collar の合併を $A_{i,t}$ とする．命題\ref{ke:global-base-parameterization}により
\[
 \sup_{y\in A_{i,t}}d_\square(f_t(y),e_i)
 \le C q_t,\qquad
 q_t=\frac{1+\log(N_{j(t)}+1)}{R_t}\longrightarrow0
\]
を満たす．従って $a_t=R_t\sqrt{\epsilon_t+q_t}$ と取れば，この全領域内の
任意の基点列について同じ成分錐への pointed GH 収束を得る．実際，
$q_t\ge R_t^{-1}$ であり，
$R_t(\epsilon_t+Cq_t)/a_t\le\max(1,C)\sqrt{\epsilon_t+q_t}\to0$ である．
この結論が使うのは全域の ambient 距離であり，core 内部の経路距離や
任意の Liouville 軌道長の一様評価は要求しない．

これは既存の大域実現定理と有限複体の局所距離からの帰結であり，
固定成分の漸近錐定理を actual 族へ直接適用する証明ではない．
特に任意に先指定した $1\ll a_t\ll R_t$ については，
頂点 fiber 上の基点では式\eqref{ke:eq:intermediate-component-scale}の誤差条件を，
一般の $A_t$ では基点のずれを含む式\eqref{ke:eq:intermediate-basepoint-budget}を
確かめる必要がある．
歪みの有効な速度は未評価なので，$t$ だけで先に決めた普遍的な中間尺度，
二つの極限の交換，又は中間尺度での tensor 収束は結論しない．

\begin{corollary}[七次以上の同じ基底における core の局所位相障害]
\label{ke:global-core-nonmatroidal}
$d\ge7$ のとき，定理\ref{ke:global-metric-realization}で選んだ $\pi_t$ の基底
$B_d(T)$ のどの core 頂点の近傍にも，二次元 matroidal local model への chart は存在しない．
従って，core 頂点を含むこの同じ基底に作用適合 matroidal atlas を入れることはできない．
これは定理\ref{trop:degree-classification}の既存の位相障害を，上の計量実現で選んだ基底に適用した帰結である．
\end{corollary}
\begin{proof}
定理\ref{ke:actual-global-fibration}の構成では，各 core 頂点 $v_i$ の近傍は，
同じ木系 $T$ から作る graph $G_i$ の cone である．命題\ref{topo:counts}より
$G_i$ は連結三価であり，
\[
 b_1(G_i)=\binom{d-2}{2}\ge10.
\]
この rank は $v_i$ における二次の局所 homology の rank でもあり，局所同相で不変である．
穿孔近傍の分岐が三価だけであることも局所位相から判定される．
一方，定理\ref{trop:cubic-classification}によれば，その分岐条件を満たす二次元
matroidal local model の link は，非連結 matroids と Euclidean factors を許しても
$b_1\le8$ を満たす．従ってそのような chart は存在しない．
ここでは作用格子や計量の収束を追加仮定せず，収縮後の $\mathcal C_d$ についての
atlas の有無も結論していない．
\end{proof}

\begin{corollary}[GH 収束と全特異 fiber 合併の収縮]
\label{ke:global-gh-and-fibers}
定理\ref{ke:global-metric-realization}の $f_t$ は全射であり，
$\epsilon_t=\operatorname{dis}(f_t)$ と置くと
\begin{align*}
 d_{\mathrm{GH}}\bigl((X_t,R_t^{-1}d_{g_t}),
                     (\mathcal C_d,d_\square)\bigr)&\le\epsilon_t/2,\\
 \sup_{c\in\mathcal C_d}
 \operatorname{diam}_{R_t^{-1}d_{g_t}}f_t^{-1}(c)&\le\epsilon_t,\\
 \sup_{b\in B_d(T)}
 \operatorname{diam}_{R_t^{-1}d_{g_t}}\pi_t^{-1}(b)&\le\epsilon_t.
\end{align*}
特に $d\ge6$ では，$p_T$ が潰す各内部木上の全ての
$\pi_t$-fibers の合併が，一様に正規化 ambient 径を零にする．
\end{corollary}

\begin{proof}
$p_T$ と $\pi_t$ は全射なので，$f_t$ の graph は両空間の correspondence である．
その distortion が $\epsilon_t$ であることから第一式を得る．
同一の $f_t$-fiber の二点を\eqref{ke:eq:global-metric-realization}に代入すると
第二式を得る．各 $\pi_t$-fiber は一つの $f_t$-fiber に含まれるので第三式も従う．
内部木の主張は命題\ref{ke:collapsed-tree-common-moment}と同じ fiber 合併であり，
命題\ref{ke:collapsed-tree-full-fiber-collapse}はその非正規化費用も与えている．
\end{proof}

\begin{corollary}[直径と径一の規格化]
\label{ke:global-diameter-limit}
$D_d=\operatorname{diam}(\mathcal C_d,d_\square)$ と置くと
\[
 R_t^{-1}\operatorname{diam}(X_t,g_t)\longrightarrow D_d,
 \qquad
 \left(X_t,\frac{d_{g_t}}{\operatorname{diam}(X_t,g_t)}\right)
 \xrightarrow{\mathrm{GH}}(\mathcal C_d,D_d^{-1}d_\square).
\]
$d\ge6$ では定理\ref{topo:diameter-six-labels}により $D_d=2\sqrt5$ である．
五次の値 $D_5=\sqrt{13}$ には，
\cite[\texttt{thm:intro-diameter-C5}]{KECubicalCurrent}の有限計量結果を用いる．
\end{corollary}

\begin{proof}
全射性と distortion 評価から
$|R_t^{-1}\operatorname{diam}(X_t,g_t)-D_d|\le\epsilon_t$ である．
$D_d>0$ なので，同じ correspondence の距離を直径で割れば第二式も従う．
\end{proof}

\paragraph{結論の範囲．}
証明は，関連 KE 原稿の未検証の大域入力
$\mathrm{(CP)}_d,\mathrm{(PBE)}_d$ 又は $\mathrm{(TC)}_5$ を仮定していない．
局所 density と全深部 triple comparison，固定 compact の収束という明記した入力から，
本稿で直接 scalar bound，source exits，基底値の比較，portal 接続を証明した．
任意に先指定した木系や critical values の実現，special Lagrangian 性，
一般次数の大域 matroidal atlas，最適収束速度は本定理の結論に含めない．
